\documentclass[10pt,a4paper]{article}

\usepackage[T1]{fontenc}
\usepackage[utf8]{inputenc}
\usepackage{textcomp}\usepackage{lmodern,amsmath,amssymb,array,tabularx,booktabs,xcolor,enumitem}
\usepackage[margin=20mm]{geometry}
\usepackage[unicode,colorlinks=true,urlcolor=black,linkcolor=black]{hyperref}
\usepackage{xcolor}
\usepackage[normalem]{ulem}
\usepackage{kotex}

\setlength{\parindent}{0pt}
\setlength{\parskip}{5pt}
\setlength{\emergencystretch}{2em}
\setlength{\tabcolsep}{4pt}
\renewcommand{\arraystretch}{1.10}
\newcommand{\hd}[1]{\par\vspace{8pt}\textbf{#1}\par\vspace{3pt}}
\newcommand{\circnum}[1]{\textcircled{\scriptsize #1}}
\newcommand{\note}{\textreferencemark\ }
\newcommand{\red}[1]{\textcolor{red}{#1}}
\newcommand{\under}[2]{\underbrace{#1}_{\text{#2}}}
\newcommand{\paper}[1]{\href{#1}{paper}}
\newcolumntype{Y}{>{\raggedright\arraybackslash}X}

%%%%%%%%%%%%%%%%%%%%%%%%%%%%%%%%%%%%%%%%%%%%%%%%%%%%%%
\begin{document}

\hd{\circnum{4} Model considering the overall market and multiple penalties}

%%%%%%%%%%%%%%%%%%%%%%%%%%%%%%%%%%%%%%%%%%%%%%%%%%%%%%
\hd{\textcolor{blue}{\textbf{[1. Construct two types of demand has been moved to pg3.]}
}}



%%%%%%%%%%%%%%%%%%%%%%%%%%%%%%%%%%%%%%%%%%%%%%%%%%%%%%
\hd{\textcolor{blue}{1}. Core model: City selection + vehicle deployment + multiple penalties}

\[
\boxed{\begin{aligned}
\min_{y,q,O,U}\quad
&\under{\sum_c F_c y_c}{Entry costs}
+\under{\sum_c K_c q_c}{Vehicle costs}\\
&+\frac{1}{|S|}\sum_s\sum_c\left[
\under{c_c^+O_c^s}{Remaining capacity}
+\under{c_c^-\textcolor{blue}{U_c^s}}{Foregone contribution margin}
\right].
\end{aligned}}
\]

%%%%%%%%%%%%%%%%%%%%%%%삭제예정%%%%%%%%%%%%%%%%%%%%%%%%
\[
{\color{red}
\boxed{\begin{aligned}
\min_{y,q,R,O,U,H,L}\quad
&\under{\sum_c F_c y_c}{Entry costs}
+\under{\sum_c K_c q_c}{Vehicle costs}\\
&+\frac{1}{|S|}\sum_s\sum_c\left[
\under{c_c^+O_c^s}{Remaining capacity}
+\under{c_{c,H}^-H_c^s}{Loss from switching to human-driven taxis}
+\under{c_{c,L}^-L_c^s}{Other unserved loss}
\right].
\end{aligned}}
}
\]
%%%%%%%%%%%%%%%%%%%%%%%%%%%%%%%%%%%%%%%%%%%%%%%%%%%%%

In $c^+,c^-$, $+,-$ distinguish excess and shortage, rather than the signs of costs. All demand, throughput, and costs use \textbf{the same evaluation period.} \note I suppose that setting the evaluation period is needed, and the model determines the city/vehicle deployment to apply during that period. 
%%%%%%%%%%%%%%%%%%%%move from pg3 to pg1%%%%%%%%%%%%%
\\

\textbf{\textcolor{blue}{2}. Constraints \textcolor{blue}{\small\textit{
[Moved from page 3 to page 1 in this revision.]
}}}

\[
\textcolor{blue}{
\boxed{\begin{aligned}
0\le q_c\le\bar q_c y_c&\qquad\forall c,\\
\sum_c q_c\le Q,&\\
\kappa_c q_c-M_c p_c^s=O_c^s-U_c^s
&\qquad\forall c,s,\\
y_c\in\{0,1\},\quad q_c\in\mathbb Z_{\ge0},&\\
O_c^s,U_c^s\ge0.&
\end{aligned}}
}
\]

{\small
\begin{tabularx}{\linewidth}{@{}>{\raggedright\arraybackslash}p{0.32\linewidth}Y@{}}
\toprule
Constraint & Meaning \\ \midrule
$0\le q_c\le\bar q_c y_c$ & Zero vehicles if the city is not entered; up to the city-specific limit if it is entered \\
$\sum_cq_c\le Q$ & Total vehicle limit \\
\textcolor{red}{\sout{$y_c\le a_c$}} & \textcolor{red}{\sout{Enter only cities where robotaxi operation is feasible at the planning stage}} \\
%%%%%%%%%%%%%%%%%%%%%%%%%%%%%%%%%%%%%%%%%%%%%
%%%%%%%%%%%%%% 삭제 예정%%%%%%%%%%%%%%%%%%%%%%
\textcolor{red}{\sout{$R_c^s+O_c^s=\kappa_c^s q_c$}}
& \textcolor{red}{\sout{Total service capacity = used capacity
+ remaining capacity (unit: trips)}} \\
\textcolor{red}{\sout{$R_c^s+U_c^s=D_c^{R,s}$}}
&\textcolor{red}{\sout{Robotaxi demand = throughput
+ unserved volume (unit: trips)}} \\
\textcolor{red}{\sout{$U_c^s=H_c^s+L_c^s$}}
&\textcolor{red}{\sout{Decompose the outcomes of unserved
demand into two types (unit: trips)}} \\
%%%%%%%%%%%%%%%%%%%%%%%%%%%%%%%%%%%%%%%%%%%%%%

\bottomrule
\end{tabularx}
}

%%%%%%%%%%%%%%%%%%%%%%%%%2pg%%%%%%%%%%%%%%%%%%%%%%%%%%%%%
\newpage
\hd{\textcolor{blue}{4}. Notation used in the model}

{\small
\begin{tabularx}{\linewidth}{@{}>{\raggedright\arraybackslash}p{0.1\linewidth}>{\raggedright\arraybackslash}p{0.37\linewidth}Y@{}}
\toprule
Symbol & Meaning & Type and unit \\ \midrule
$C$ & Set of potential candidate cities & Input set; cities \\
$S$ & Set of demand scenarios & Input set; scenarios \\
$y_c$ & 1 if city $c$ is entered, else 0 & Decision variable; binary (0 or 1) \\
$q_c$ & Number of vehicles deployed in city $c$ & Decision variable; vehicles (integer) \\
$\kappa_c$ & Number of trips one vehicle can serve during the period & Input parameter; trips/vehicle over the evaluation period. \newline \red{\textbf{Operational data or assumptions or fixed?}}  \\
$Q$ & Total number of vehicles available for deployment & Input parameter; vehicles \\
$\bar q_c$ & Upper bound on vehicle deployment in each city & Input parameter; vehicles \\

%%%%%%%%%%%%%% 비용 항목 %%%%%%%%%%%%%%%%%
$F_c$
& Fixed cost of city entry and basic infrastructure (도시c에 진입하는 비용)
& Input parameter; currency over the evaluation period
\\

$K_c$
& Per-vehicle preparation, deployment, and holding cost
over the evaluation period(감가상각비, 보험료, 고정주차 및 관리 비용)
& Input parameter; currency/vehicle over the evaluation period
\\

$c_c^+$
& Additional loss per unit of unused service capacity (남는 용량 때문에 추가로 발생하는 손실 / 유휴 서비스 용량에 대한 손실 평가계수)
& Input parameter; currency/trip of unused capacity
\\

$c_c^-$
& Foregone contribution margin per unit of unserved
potential demand
& Input parameter; currency/unserved trip
\\

\textcolor{blue}{$M_c$} & Number of \textbf{potential ride-hailing usage opportunities} in city $c$ & Input parameter; usage opportunities over the evaluation period. \newline \textcolor{blue}{data or assumption} \\
%%%%%%%%%%%%%%%%%%%%%%%%%%%%%%%%%%%%
%%%%%%%%%%%%%삭제 예정%%%%%%%%%%%%%%%
\textcolor{red}{\sout{$a_c$}}
& \textcolor{red}{\sout{1 if operation in the city is feasible
at the planning stage, else 0}}
& \textcolor{red}{\sout{Assessed using permit, ODD, and facility
data. Separate assumptions for unverified conditions}} \\

\textcolor{red}{\sout{$B$}}
& \textcolor{red}{\sout{Upper bound on the number of cities
that can be entered simultaneously}}
& \textcolor{red}{\sout{Company input or an experimental value
\textbf{only when an actual restriction exists}}} \\
%%%%%%%%%%%%%%%%%%%%%%%%%%%%%%%%%%%%
\textcolor{red}{\sout{\textcolor{blue}{$M_c$}}} & \textcolor{red}{\sout{Number of \textbf{potential ride-hailing usage opportunities} in city $c$, scenario $s$}} & \textcolor{red}{\sout{(Input) Need to be Estimated}} \\
$p_c^s$ & Probability of choosing our company's robotaxi for that usage opportunity & Scenario input; dimensionless probability in $[0,1]$ \\
$D_c^s$ & Robotaxi choice demand & Derived input; trips over the evaluation period. \newline \textcolor{blue}{$D_c^s= M_cp_c^s$} \\
$D_c^{N,s}$ & Demand not choosing robotaxis (human-driven taxis, other modes, giving up the trip, etc.) & Derived input; trips over the evaluation period. \newline \textcolor{blue}{$D_c^{N,s}= M_c(1-p_c^s)$} \\
%%%%%%%%%%%%%%%%%%%삭제 예정%%%%%%%%%%%%%%%%%%
\textcolor{red}{\sout{$\tau_c^s$}}
& \textcolor{red}{\sout{Proportion of unserved customers
who switch to human-driven taxis}}
& \textcolor{red}{\sout{\textbf{estimated by data or fixed?}}} \\
%%%%%%%%%%%%%%%%%%%%%%%%%%%%%%%%%%%%%%%%%%%%
%%%%%%%%%%%%%%%%%%삭제 예정%%%%%%%%%%%%%%%%%%
\textcolor{red}{\sout{$R_c^s$}}
& \textcolor{red}{\sout{Actual robotaxi throughput}}
& \textcolor{red}{\sout{Calculated by the model based on
the deployment outcome}} \\
%%%%%%%%%%%%%%%%%%%%%%%%%%%%%%%%%%%%%%%%%%%

$O_c^s$ & Unused service capacity
& Outcome variable; trips over the evaluation period. \newline \textcolor{blue}{$O_c^s =(\kappa_c q_c-D_c^s)^+$} \\

$U_c^s$ & Unserved robotaxi demand
& Outcome variable; trips over the evaluation period. \newline \textcolor{blue}{$U_c^s=(D_c^s-\kappa_c q_c)^+$} \\

%%%%%%%%%%% 삭제 예정 %%%%%%%%%%%%%%%%%%%%%%
\textcolor{red}{\sout{$H_c^s$}}
& \textcolor{red}{\sout{Demand that switches to human-driven
taxis and is served after being unserved}}
& \textcolor{red}{\sout{Calculated using the switching
condition below $(H_c^s=\tau_c^s U_c^s)$}} \\

\textcolor{red}{\sout{$L_c^s$}}
& \textcolor{red}{\sout{Remaining unserved demand}}
& \textcolor{red}{\sout{Groups other modes, giving up the
trip, etc. $(L_c^s=(1-\tau_c^s)U_c^s)$}} \\
\bottomrule
\end{tabularx}
}

\note Demand-related variables consist of the inputs $p_c^s,M_c$ and outcome variables calculated from scenario demand and deployment in the model.
\note Outcome variables are not quantities to be predicted by collecting data for each one separately. They are outcomes calculated from demand and deployment.

%%%%%%%%%%%%%%%%%%%%Move from pg1 to pg3%%%%%%%%%%%%%%%%%%%%%%%%%
\newpage
\hd{\textcolor{blue}{5. Explanations}}

\textbf{A. Construct two types of demand \textcolor{blue}{\small\textit{
[Moved from page 3 to page 1 in this revision.]
}}}

\[
D_c^s = \textcolor{blue}{M_c} p_c^s,\qquad 
D_c^{N,s}= \textcolor{blue}{M_c} (1-p_c^s)
\]

{\small
\begin{tabularx}{\linewidth}{@{}>{\raggedright\arraybackslash}p{0.12\linewidth}>{\raggedright\arraybackslash}p{0.42\linewidth}Y@{}}
\toprule
Symbol & Meaning & Explanation \\ \midrule
\textcolor{blue}{$M_c$} & Number of \textbf{potential ride-hailing usage opportunities} in city $c$, \textcolor{red}{\sout{scenario $s$}} & Robotaxi demand + human-driven taxi demand + other demand \\
$p_c^s$ & Probability of choosing robotaxi for that usage opportunity & Need to be specified with reference service conditions such as fares and waiting times \\
$D_c^s$ & Robotaxi choice demand & Demand \textbf{before} applying vehicle shortages \\
$D_c^{N,s}$ & Demand not choosing robotaxis & Human-driven taxis, other modes, giving up the trip, etc. \\
\bottomrule
\end{tabularx}
}

%%%%%%%%%%%%%%%%%%%%%%%%%%%%%%%%%%%%%%%%%%%%%%%%%%%%%%%%%%%%%
\hd{\textcolor{red}{\sout{4. Common constraints}}}

\textbf{B. Constraints}\\
1) Once demand $D_c^s$ and service capacity $\kappa_c q_c$ are determined, the following two values are calculated based on the deployment outcome; they are not variables that need to be investigated separately. $(x)^+=\max\{x,0\}$. Both values are measured in trips.

\[
\begin{aligned}
%%%%%%%%%% 삭제 필요 %%%%%%%%%%%%%%%%%
\textcolor{red}{\mbox{\sout{\(R_c^s\)}}}
&\textcolor{red}{\mbox{\sout{\(
  \displaystyle = \min\{D_c^{R,s},\kappa_c q_c\}
\)}}}
&&\textcolor{red}{\sout{Actual throughput}},\\
%%%%%%%%%%%%%%%%%%%%%%%%%%%%%%%%%%%%%%
O_c^s&=(\kappa_c q_c-D_c^s)^+&&\text{Remaining service capacity},\\
U_c^s&=(D_c^s-\kappa_c q_c)^+&&\text{Unserved potential demand}.
\end{aligned}
\]\\
%%%%%%%%%% 삭제 필요 %%%%%%%%%%%%%%%%%%
\begingroup
\color{red}
\sout{
2) For human fallback, I set the switching proportion $\tau_c^s$ as an input.
}

\[
\mbox{\sout{\(
\displaystyle H_c^s=\tau_c^s U_c^s,\qquad L_c^s=(1-\tau_c^s)U_c^s
\)}} 
\]

\sout{
$\tau_c^s\in[0,1]$ is the proportion of unserved demand that actually switches to human-driven taxis and is served.
}
\par
\sout{
Then, the cost terms gives
}

\[
\mbox{\sout{\(
\displaystyle
c_{c,H}^-H_c^s+c_{c,L}^-L_c^s
=\underbrace{\left[\tau_c^s c_{c,H}^-+(1-\tau_c^s)c_{c,L}^-\right]}_{\text{Effective shortage cost}}U_c^s.
\)}}
\] \\

\sout{
This is in the case when $\tau_c^s$ is fixed (a constant irrelevant of scenario, city selection and so on).
}

\par
\endgroup
%%%%%%%%%%%%%%%%%%%%%%%%%%%%%%%%%%%%%%%%%%%%%%%%%%%%

\textcolor{blue}{\textbf{[B. Constraints has been moved to page 1.]}
}

%%%%%%%%%%%%%%%%%%%%%%%%%%%%%%%%%%%%%%%%%%%%%%%%%%%%%
\hd{\textcolor{red}{\sout{\note How to set inputs?}}}
\textcolor{red}{\sout{I wrote about my thoughts of how I might set the inputs.}}

{\small
\begin{tabularx}{\linewidth}{@{}>{\raggedright\arraybackslash}p{0.1\linewidth}>{\raggedright\arraybackslash}p{0.25\linewidth}>{\raggedright\arraybackslash}p{0.29\linewidth}Y@{}}
\toprule
\textcolor{red}{\sout{Input}} & \textcolor{red}{\sout{How to set it}} & \textcolor{red}{\sout{Use of previously collected data}} & \textcolor{red}{\sout{Data still required}} \\ \midrule
\textcolor{red}{\sout{\textcolor{blue}{$M_c$}}} & \textcolor{red}{\textbf{\sout{Data estimation + market-definition assumptions}}} & \textcolor{red}{\sout{Population, employment, existing ride-hailing, airports, commercial activity}} & \textcolor{red}{\sout{Denominator of the potential travel market. Completed trips alone are insufficient}} \\
\textcolor{red}{\sout{$p_c^{R,s}$}} & \textcolor{red}{\textbf{\sout{Estimation from choice data;}}\sout{ assumptions if unavailable}} & \textcolor{red}{\sout{Candidate explanatory variables such as income, car-free households, and transport accessibility}} & \textcolor{red}{\sout{Robotaxi choice/non-choice, offered fares, waiting times, and alternatives}} \\

\textcolor{red}{\sout{\textcolor{blue}{$\kappa_c^s$}}} & \textcolor{red}{\textbf{\sout{Estimation from operational data;}}\sout{ assumptions if unavailable}} & \textcolor{red}{\sout{Congestion, roads, weather, charging conditions}} & \textcolor{red}{\sout{Driving, pickup, repositioning, charging, and maintenance times}} \\

%%%%%%%%%%%%%%% 삭제 예정 %%%%%%%%%%%%%%%%%%%
\textcolor{red}{\sout{$\tau_c^s$}}
&{\color{red}
{\bfseries\sout{Estimation from switching data;}}
\sout{assumptions if unavailable}}
&\textcolor{red}{\sout{Cannot be estimated directly
from the current data}}
&\textcolor{red}{\sout{Actual behavior after robotaxi
demand goes unserved}} \\
\textcolor{red}{\sout{$a_c$}}
&{\color{red}\bfseries
\sout{Assessment from data + assumptions about launch conditions}}
&\textcolor{red}{\sout{Permit, ODD, and base conditions}}
&\textcolor{red}{\sout{Conditions for the company
at the planning stage}} \\
%%%%%%%%%%%%%%%%%%%%%%%%%%%%%%%%%%%%%%%%%%%%

\textcolor{red}{\sout{$\bar q_c$}} & \textcolor{red}{\textbf{\sout{Facility/operational data or company input}}} & \textcolor{red}{\sout{Charging, parking, and maintenance conditions}} & \textcolor{red}{\sout{Actual scale that can be secured. $Q$ if there is no separate upper bound}} \\
\textcolor{red}{\sout{$Q$}} & \textcolor{red}{\textbf{\sout{Company input or experimental value}}} & \textcolor{red}{\sout{Not estimated from demand data}} & \textcolor{red}{\sout{Number of available vehicles}} \\
%%%%%%%%%%%%%%%%%%%삭제예정%%%%%%%%%%%%%%%%
\textcolor{red}{\sout{$F_c$}} & \textcolor{red}{\textbf{\sout{Additional cost data or assumptions}}} & \textcolor{red}{\sout{Some base and land costs}} & \textcolor{red}{\sout{Initial construction and entry preparation costs}} \\
\textcolor{red}{\sout{$K_c$}} & \textcolor{red}{\textbf{\sout{Additional cost data or assumptions}}} & \textcolor{red}{\sout{Cannot be determined..}} & \textcolor{red}{\sout{Vehicle deployment and holding costs for the evaluation period}} \\
\textcolor{red}{\sout{$c_c^+$}} & \textcolor{red}{\textbf{\sout{Company loss data or assumptions}}} & \textcolor{red}{\sout{Cannot be determined..}} & \textcolor{red}{\sout{Burden of remaining capacity. No double counting with $K_c$}} \\
\textcolor{red}{\sout{$c_{c,H}^-,c_{c,L}^-$}} & \textcolor{red}{\textbf{\sout{Company loss data or assumptions}}} & \textcolor{red}{\sout{Cannot be determined..}} & \textcolor{red}{\sout{Differences in opportunity loss and customer churn by switching type}} \\
%%%%%%%%%%%%%%%%%%%%%%%%%%%%%%%%%%%%%%%%%%
\bottomrule
\end{tabularx}
}

\textcolor{red}{\sout{\note{$y_c,q_c$ are not included; values determined by the model}}}

%%%%%%%%%%%%%%%%%%%%%%%%%%%%%%%%%%%%%%%%%%%%%%%%%%%%%%
\newpage
\hd{5. Driving factors}
\hd{A. Driving factors of $M$}
The papers below mainly explain \textbf{observed ride-hailing usage}, so they do not include the `unrealized demand' included in the definition of $M$, making direct estimation impossible. However, I suppose the demand patterns provided by these papers lead to constructing `potential travel opportunities $M$' through modeling.

{\small
\begin{tabularx}{\linewidth}{@{}>{\raggedright\arraybackslash}p{0.16\linewidth}>{\raggedright\arraybackslash}p{0.32\linewidth}>{\raggedright\arraybackslash}p{0.3\linewidth}Y@{}}
\toprule
Variable & Supporting paper & Scope established in the paper & Measurement data \\ \midrule
\textbf{Population and population density} & Elkhouly and Alhadidi, \textit{Ride-sharing Determinants} (\paper{https://arxiv.org/abs/2406.11590}); Zhang et al., \textit{Examining spatial heterogeneity\ldots} (\paper{https://arxiv.org/abs/2209.07980}) & Relationship of population and local density with conventional ride-hailing usage. Density effects may vary by area & \href{https://www.census.gov/programs-surveys/acs}{ACS} \\
\textbf{Employment and economic activity} & Elkhouly and Alhadidi (\paper{https://arxiv.org/abs/2406.11590}) & Uses the economically active population and employment environment to explain usage & \href{https://lehd.ces.census.gov/data/}{LEHD/LODES}, ACS \\
\textbf{Income} & Elkhouly and Alhadidi (\paper{https://arxiv.org/abs/2406.11590}) & Analyzes the relationship between income and usage in Chicago & \href{https://data.census.gov/table/ACSDT5Y2024.B19013}{ACS B19013} \\
\textbf{Share of car-free households} & Elkhouly and Alhadidi (\paper{https://arxiv.org/abs/2406.11590}) & Analyzes the relationship between households without cars and usage & \href{https://data.census.gov/table/ACSDT1Y2022.B08201}{ACS B08201} \\
\textbf{Land-use mix, accessibility, and commercial activity} & Zhang et al., \textit{Examining spatial heterogeneity\ldots} (\paper{https://arxiv.org/abs/2209.07980}) & Nonlinear and spatial differences in the urban environment and ridesourcing usage & \href{https://www.epa.gov/smartgrowth/smart-location-mapping}{EPA SLD} \\
\textbf{Restaurant and leisure-facility density} & Loa, Lee, and Circella, \textit{Just another ride-hailing?} (\paper{https://link.springer.com/article/10.1007/s11116-026-10751-9}) & Relationship with \textbf{robotaxi trip counts in SF}. Does not separate the effects of $M$ and $p$ & \href{https://www.census.gov/programs-surveys/cbp.html}{Census CBP}, OpenStreetMap \\

\textbf{Public transport and walking environment} & Zhang et al. (\paper{https://arxiv.org/abs/2209.07980}); Cats et al., \textit{Beyond the dichotomy\ldots} (\paper{https://doi.org/10.1371/journal.pone.0262496}) & Analyzes the local environment and travel-time competitiveness and substitutability with public transport. Not a uniform increase/decrease effect & EPA SLD, \href{https://gtfs.org/documentation/overview/}{GTFS} \\
\textbf{Airport-related market size} & Zhang et al. (\paper{https://arxiv.org/abs/2209.07980}) & Analyzes airport travel as a separate spatial category. \textbf{Does not estimate the direct effect of airport passenger counts} & \href{https://www.faa.gov/airports/planning_capacity/passenger_allcargo_stats/passenger}{FAA passenger data}. Used as a proxy for market size \\
\textbf{Time, day of week, and season} & Elkhouly and Alhadidi (\paper{https://arxiv.org/abs/2406.11590}) & Analyzes differences in usage over time. Specific variables such as holidays are selected through separate design & Dates and times in trip data \\
\textbf{Weather} & Additional: Liu, Jiang, and Chen, \textit{Quantifying the impact of weather on ride-hailing ridership: Evidence from Haikou, China} (2021, \paper{https://doi.org/10.1016/j.tbs.2021.04.002}) & Analyzes the relationship between weather and ride-hailing usage. Effects in U.S. cities require revalidation & \href{https://www.ncei.noaa.gov/cdo-web/}{NOAA CDO} \\
\bottomrule
\end{tabularx}
}

%The determinants of $M$ are city characteristics (original proposal) + temporal/seasonal context (situation) + market activity (outcome: past Uber/Lyft usage) Past usage is a predictor that can be used separately. However, this does not explain ``why demand occurs''; it can be classified as \textbf{outcome data showing the size (scale) of demand.}

%%%%%%%%%%%%%%%%%%%%%%%%%%%%%%%%%%%%%%%%
\newpage
\hd{B. Driving factors of $p$}

{\small
\begin{tabularx}{\linewidth}{@{}>{\raggedright\arraybackslash}p{0.16\linewidth}>{\raggedright\arraybackslash}p{0.32\linewidth}>{\raggedright\arraybackslash}p{0.3\linewidth}Y@{}}
\toprule
Variable & Supporting paper & Scope of application & Required measurement data \\ \midrule
\textbf{Relative fare} & Additional: Yin and Cherchi, \textit{Willingness to pay for automated taxis\ldots} (2024; online 2022, \paper{https://doi.org/10.1007/s11116-022-10319-3}) & Choice experiment between conventional taxis and automated taxis & Offered fares for both + choice outcomes \\
\textbf{Relative waiting time} & Yin and Cherchi (\paper{https://doi.org/10.1007/s11116-022-10319-3}) & Analyzes choices under different service conditions & Offered waiting times for both + choice outcomes \\
\textbf{In-vehicle travel time} & Yin and Cherchi (\paper{https://doi.org/10.1007/s11116-022-10319-3}) & Analyzes choice trade-offs between fares and time & Offered travel times for both \\
\textbf{Trust, price value, and usage habits} & Additional: Wei et al., \textit{Using the Extended Unified Theory\ldots} (2024, \paper{https://www.nature.com/articles/s41599-024-03271-3}) & Study of robotaxi acceptance intention and behavior. \textbf{Survey scales cannot be interpreted directly as choice probabilities} & Individual attitude and habit surveys \\
\textbf{Age and robotaxi-related knowledge} & Yin and Cherchi (\paper{https://doi.org/10.1007/s11116-022-10319-3}) & Choice differences by individual characteristics. Fixing the sign, such as ``younger people always prefer it, is inappropriate & Link individual characteristics to choice data \\
\textbf{User reviews and in-vehicle support features} & Yin and Cherchi (\paper{https://doi.org/10.1007/s11116-022-10319-3}) & Analyzes customer ratings and features such as destination changes and communication with an operator as choice attributes & Service attributes in the choice experiment \\
\textbf{Experience using robotaxis} & Additional: Yin, Cherchi, and Nettle, \textit{Living in a Ci ty Where Automated Taxis are Operating and Using Them\ldots} (2023, \paper{https://doi.org/10.1177/03611981231155181}) & Analyzes preference differences by whether service operates in the residential area and by individual usage experience & Individual usage experience \\

\textbf{Competitiveness of alternatives such as public transport} & Cats et al., \textit{Beyond the dichotomy\ldots} (\paper{https://doi.org/10.1371/journal.pone.0262496}) & Evidence from \textbf{conventional ride-hailing}. Not a robotaxi-specific preference effect & Availability, fares, and times of alternatives \\
\bottomrule
\end{tabularx}
}


%%%%%%%%%%%%%%%%%%%%%%%%%%%%%%%%%%%%%%%%%%%%%%%%%%%%%%
\textcolor{red}{\sout{\hd{6. Things needed to be clarified} \\
1. how to set scenarios?($M$,$p$,$\tau$, $\kappa$)\\ -> Kannan's paper? \\
2. What can be estimated from the existing research materials?\\
3. What additional data are needed?\\
4. How should assumptions be made for symbols for which data are unavailable?\\
5. Can a solution be obtained for the objective-function formulation?\\
6. If data cannot be obtained for the remaining symbols other than the outcome variables that need to be estimated, how should they be set? \\
7. How to choose which factors to use for model?
}}
\end{document}
